\begin{textsample}{POS Dim 1 – 2012 – Score 20.00 – 2012/t002268.txt}  \label{ex:f1_pos_037}
\textbf{The} job of uncovering \textbf{the} global food waste scandal started for me when I was 15 years old.
I bought some pigs.
I was living in Sussex.
And I started to feed them in \textbf{the} most traditional and environmentally \textbf{friendly} way.
I went to my school kitchen, and I said,"Give me \textbf{the} scraps that my school friends have turned their noses up at."I went to \textbf{the} local baker and took their stale bread.
I went to \textbf{the} local greengrocer, and I went to a farmer who was throwing away potatoes because they were \textbf{the} wrong shape or size for \textbf{supermarkets}.
This was great.
My pigs turned that food waste into delicious pork.
I sold that pork to my school friends ' parents, and I made a good \textbf{pocket} money addition to my teenage allowance.
But I noticed that most of \textbf{the} food that I was giving my pigs was in fact fit for human \textbf{consumption}, and that I was only scratching \textbf{the} \textbf{surface}, and that right \textbf{the} way up \textbf{the} food supply chain, in \textbf{supermarkets}, greengrocers, bakers, in our homes, in factories and farms, we were hemorrhaging out food.
Supermarkets didn ' t even want to talk to me about how much food they were wasting.
I ' d been round \textbf{the} back.
I ' d seen bins full of food being locked and then trucked off to landfill sites, and I thought, surely there is something more sensible to do with food than waste it.
One morning, when I was feeding my pigs, I noticed a particularly tasty-looking sun-dried tomato loaf that used to crop up from time to time.
I grabbed hold of it, sat down, and ate my breakfast with my pigs.
That was \textbf{the} first act of what I later learned to call freeganism, really an exhibition of \textbf{the} injustice of food waste, and \textbf{the} provision of \textbf{the} solution to food waste, which is simply to sit down and eat food, rather than throwing it away.
That became, as it were, a way of confronting large businesses in \textbf{the} business of wasting food, and exposing, most importantly, to \textbf{the} public, that when we ' re talking about food being thrown away, we ' re not talking about rotten stuff, we ' re not talking about stuff that ' s beyond \textbf{the} pale.
We ' re talking about good, fresh food that is being wasted on a colossal scale.
Eventually, I set about writing my book, really to demonstrate \textbf{the} extent of this problem on a global scale.
What this shows is a nation-by-nation breakdown of \textbf{the} likely level of food waste in each country in \textbf{the} world.
Unfortunately, empirical data, good, hard stats, don ' t exist, and therefore to prove my point, I first of all had to find some proxy way of uncovering how much food was being wasted.
So I took \textbf{the} food supply of every single country and I compared it to what was actually likely to be being consumed in each country.
That ' s based on diet intake surveys, it ' s based on levels of obesity, it ' s based on a range of factors that gives you an approximate guess as to how much food is actually going into people ' s mouths.
That black line in \textbf{the} middle of that table is \textbf{the} likely level of \textbf{consumption} with an allowance for certain levels of inevitable waste.
There will always be waste.
I ' m not that unrealistic that I think we can live in a waste-free world.
But that black line shows what a food supply should be in a country if they allow for a good, stable, secure, nutritional diet for every person in that country.
Any dot above that line, and you ' ll quickly notice that that includes most countries in \textbf{the} world, represents unnecessary surplus, and is likely to reflect levels of waste in each country.
As a country gets richer, it invests more and more in getting more and more surplus into its shops and restaurants, and as you can see, most European and North American countries fall between 150 and 200 percent of \textbf{the} nutritional requirements of their populations.
So a country like America has twice as much food on its shop shelves and in its restaurants than is actually required to feed \textbf{the} American people.
But \textbf{the} thing that really struck me, when I plotted all this data, and it was a lot of numbers, was that you can see how it levels off.
Countries rapidly shoot towards that 150 mark, and then they level off, and they don ' t really go on rising as you might expect.
So I decided to unpack that data a little bit further to see if that was true or false.
And that ' s what I came up with.
If you include not just \textbf{the} food that ends up in shops and restaurants, but also \textbf{the} food that people feed to livestock, \textbf{the} maize, \textbf{the} soy, \textbf{the} wheat, that humans could eat but choose to fatten livestock instead to produce increasing amounts of meat and dairy products, what you find is that most rich countries have between three and four times \textbf{the} amount of food that their population needs to feed itself.
A country like America has four times \textbf{the} amount of food that it needs.
When people talk about \textbf{the} need to increase global food production to feed those nine billion people that are expected on \textbf{the} \textbf{planet} by 2050, I always think of these graphs. \textbf{The} fact is, we have an enormous buffer in rich countries between ourselves and hunger.
We ' ve never had such gargantuan surpluses before.
In many ways, this is a great success story of human civilization, of \textbf{the} agricultural surpluses that we set out to achieve 12, 000 years ago.
It is a success story.
It has been a success story.
But what we have to recognize now is that we are reaching \textbf{the} ecological limits that our \textbf{planet} can bear, and when we chop down forests, as we are every day, to grow more and more food, when we extract water from depleting water reserves, when we emit \textbf{fossil} fuel \textbf{emissions} in \textbf{the} quest to grow more and more food, and then we throw away so much of it, we have to think about what we can start saving.
And yesterday, I went to one of \textbf{the} local \textbf{supermarkets} that I often visit to inspect, if you like, what they ' re throwing away.
I found quite a few packets of biscuits amongst all \textbf{the} fruit and vegetables and everything else that was in there.
And I thought, well this could serve as a symbol for today.
So I want you to imagine that these nine biscuits that I found in \textbf{the} bin represent \textbf{the} global food supply, \textbf{okay}?
We start out with nine.
That ' s what ' s in fields around \textbf{the} world every single year. \textbf{The} first biscuit we ' re going to lose before we even leave \textbf{the} farm.
That ' s a problem primarily associated with developing work \textbf{agriculture}, whether it ' s a lack of infrastructure, refrigeration, pasteurization, grain stores, even basic fruit crates, which means that food goes to waste before it even leaves \textbf{the} fields. \textbf{The} next three biscuits are \textbf{the} foods that we decide to feed to livestock, \textbf{the} maize, \textbf{the} wheat and \textbf{the} soya.
Unfortunately, our beasts are inefficient animals, and they turn \textbf{two-thirds} of that into feces and heat, so we ' ve lost those two, and we ' ve only kept this one in meat and dairy products.
Two more we ' re going to throw away directly into bins.
This is what most of us think of when we think of food waste, what ends up in \textbf{the} garbage, what ends up in \textbf{supermarket} bins, what ends up in restaurant bins.
We ' ve lost another two, and we ' ve left ourselves with just four biscuits to feed on.
That is not a superlatively efficient use of global resources, especially when you think of \textbf{the} billion hungry people that exist already in \textbf{the} world.
Having gone through \textbf{the} data, I then needed to demonstrate where that food ends up.
Where does it end up?
We ' re used to seeing \textbf{the} stuff on our \textbf{plates}, but what about all \textbf{the} stuff that goes missing in between? \textbf{Supermarkets} are an easy place to start.
This is \textbf{the} result of my hobby, which is unofficial bin inspections.
Strange you might think, but if we could rely on corporations to tell us what they were doing in \textbf{the} back of their stores, we wouldn ' t need to go sneaking around \textbf{the} back, opening up bins and having a look at what ' s inside.
But this is what you can see more or less on every street corner in Britain, in Europe, in North America.
It represents a colossal waste of food, but what I discovered whilst I was writing my book was that this very evident abundance of waste was actually \textbf{the} tip of \textbf{the} iceberg.
When you start going up \textbf{the} supply chain, you find where \textbf{the} real food waste is happening on a gargantuan scale.
Can I have a show of hands if you have a loaf of sliced bread in your house?
Who lives in a household where that crust — that slice at \textbf{the} first and last end of each loaf — who lives in a household where it does get eaten? \textbf{Okay}, most people, not everyone, but most people, and this is, I ' m glad to say, what I see across \textbf{the} world, and yet has anyone seen a \textbf{supermarket} or sandwich shop anywhere in \textbf{the} world that serves sandwiches with crusts on it?
I certainly haven ' t.
So I kept on thinking, where do those crusts go?
This is \textbf{the} answer, unfortunately : 13, 000 slices of fresh bread coming out of this one single factory every single day, day-fresh bread.
In \textbf{the} same year that I visited this factory, I went to Pakistan, where people in 2008 were going hungry as a result of a squeeze on global food supplies.
We contribute to that squeeze by depositing food in bins here in Britain and elsewhere in \textbf{the} world.
We take food off \textbf{the} market shelves that hungry people depend on.
Go one step up, and you get to farmers, who throw away sometimes a third or even more of their harvest because of cosmetic standards.
This farmer, for example, has invested 16, 000 \textbf{pounds} in growing spinach, not one leaf of which he harvested, because there was a little bit of grass growing in amongst it.
Potatoes that are cosmetically imperfect, all going for pigs.
Parsnips that are too small for \textbf{supermarket} specifications, tomatoes in Tenerife, oranges in Florida, bananas in Ecuador, where I visited last year, all being \textbf{discarded}.
This is one day ' s waste from one banana plantation in Ecuador.
All being \textbf{discarded}, perfectly edible, because they ' re \textbf{the} wrong shape or size.
If we do that to fruit and vegetables, you bet we can do it to animals too.
Liver, lungs, heads, tails, kidneys, testicles, all of these things which are traditional, delicious and nutritious parts of our gastronomy go to waste.
Offal \textbf{consumption} has halved in Britain and America in \textbf{the} last 30 years.
As a result, this stuff gets fed to dogs at best, or is incinerated.
This man, in Kashgar, Xinjiang province, in Western China, is serving up his national dish.
It ' s called sheep ' s organs.
It ' s delicious, it ' s nutritious, and as I learned when I went to Kashgar, it symbolizes their taboo against food waste.
I was sitting in a roadside cafe.
A chef came to talk to me, I finished my bowl, and halfway through \textbf{the} conversation, he stopped talking and he started frowning into my bowl.
I thought,"My goodness, what taboo have I broken?
How have I insulted my host?"He pointed at three grains of rice at \textbf{the} bottom of my bowl, and he said,"Clean."I thought,"My God, you know, I go around \textbf{the} world telling people to stop wasting food.
This guy has thrashed me at my own game."But it gave me faith.
It gave me faith that we, \textbf{the} people, do have \textbf{the} power to stop this tragic waste of resources if we regard it as socially unacceptable to waste food on a colossal scale, if we make noise about it, tell corporations about it, tell governments we want to see an end to food waste, we do have \textbf{the} power to bring about that change. \textbf{Fish}, 40 to 60 percent of European \textbf{fish} are \textbf{discarded} at \textbf{sea}, they don ' t even get landed.
In our homes, we ' ve lost touch with food.
This is an experiment I did on three lettuces.
Who keeps lettuces in their fridge?
Most people. \textbf{The} one on \textbf{the} \textbf{left} was kept in a fridge for 10 days. \textbf{The} one in \textbf{the} middle, on my kitchen table.
Not much difference. \textbf{The} one on \textbf{the} right I treated like cut flowers.
It ' s a living organism, cut \textbf{the} slice off, stuck it in a vase of water, it was all right for another two weeks after this.
Some food waste, as I said at \textbf{the} beginning, will inevitably arise, so \textbf{the} question is, what is \textbf{the} best thing to do with it?
I answered that question when I was 15.
In fact, humans answered that question 6, 000 years ago : We domesticated pigs to turn food waste back into food.
And yet, in Europe, that practice has become illegal since 2001 as a result of \textbf{the} foot-and-mouth outbreak.
It ' s unscientific.
It ' s unnecessary.
If you cook food for pigs, just as if you cook food for humans, it is \textbf{rendered} safe.
It ' s also a massive saving of resources.
At \textbf{the} moment, Europe depends on importing millions of tons of soy from South America, where its production contributes to global warming, to deforestation, to biodiversity loss, to feed livestock here in Europe.
At \textbf{the} same time we throw away millions of tons of food waste which we could and should be feeding them.
If we did that, and fed it to pigs, we would save that amount of \textbf{carbon}.
If we feed our food waste which is \textbf{the} current government favorite way of getting rid of food waste, to anaerobic digestion, which turns food waste into gas to produce electricity, you save a paltry 448 kilograms of \textbf{carbon} dioxide per ton of food waste.
It ' s much better to feed it to pigs.
We knew that during \textbf{the} war.
A silver lining : It has kicked off globally, \textbf{the} quest to tackle food waste.
Feeding \textbf{the} 5, 000 is an event I first organized in 2009.
We fed 5, 000 people all on food that otherwise would have been wasted.
Since then, it ' s happened again in London, it ' s happening internationally, and across \textbf{the} country.
It ' s a way of organizations coming together to celebrate food, to say \textbf{the} best thing to do with food is to eat and \textbf{enjoy} it, and to stop wasting it.
For \textbf{the} sake of \textbf{the} \textbf{planet} we live on, for \textbf{the} sake of our children, for \textbf{the} sake of all \textbf{the} other organisms that share our \textbf{planet} with us, we are a terrestrial animal, and we depend on our land for food.
At \textbf{the} moment, we are trashing our land to grow food that no one eats.
Stop wasting food.
Thank you very much.

% matched lemmas: agriculture, carbon, consumption, discard, emission, enjoy, fish, fossil, friendly, left, okay, planet, plate, pocket, pound, render, sea, supermarket, surface, the, two-third
\end{textsample}
